%!TEX root = ../dokumentation.tex

\chapter{Implementierung der Netzwerkinfrastruktur}

\section{Lokale IoT-Datenübertragung}
Wie der ESP32 sich ins WLAN einwählt, die TCP-Verbindung aufbaut und die Sensordaten per MQTT an den Port 1883 des Raspberry Pis sendet.

\section{Interne Systemkommunikation (Docker Routing)}
Die Umsetzung der Theorie aus 2.5. Grafana und Telegraf finden die InfluxDB einfach über den Container-Namen (z.B. http://influxdb:8086). Daten verlassen hierfür nie das virtuelle docker0-Netz. Docker netz mit einer .yaml DAtei strukturiert, dabei benötigte berechigungen und APU-Keys Übergeben (mehr dazu im Sicherheitsteil).

\section{Externe Datenübertragung (Cloudflare Tunnel)}
Die praktische Umsetzung von 2.6. Installation des Cloudflared-Dienstes auf dem Pi, Verknüpfung mit der gekauften Domain und das automatische HTTPS-Zertifikat für den sicheren Browser-Zugriff.

\section{Raspberry Pi als WLAN-Access Point}
Das ganze Problem mir nicht genug Strom vom Netzteil. Warum? IPhone Hotspot auf 5 Verbindungen begrenzt und es ist näher an einer reellen IoT Anwendung, in einem Auto.

